\documentclass{journal}

\usepackage{graphicx}
\graphicspath{{../../drills}}


\usepackage{amsmath}
\usepackage[nopatch]{microtype}
\usepackage{booktabs}
\usepackage{hyperref}

\usepackage[left=2.45cm,top=3cm,right=2.45cm,bottom=2.5cm,bindingoffset=0cm]{geometry}

\title{Intro to Powerjams}

\author{}

\begin{document}
\maketitle
\noindent 

This training session focuses on basic power jam strategies.

\section*{Overview}

\subsection*{Pre-Session: Off Skates Warm-up 5-10 Minutes}
Lightly making sure that the joints still work.

Focus is {\bf not} on building strength, cardio or endurance, just on limbering up. Do 5 to 10 reps of each, on each side. 

\begin{itemize}
\item Toe taps + one leg light squats
\item Arm, shoulder and neck stretches
\item Opening and closing hip rotations
\item Knee raise to superman 
\item Cross chest knee raises
\item Butt kicks, high knees, side shuffles
\end{itemize}


\subsection*{Warm Up - 10 Minutes}

This session will involve some crossovers skating, particularly backwards, so this is included in the warm-up.

\begin{itemize}
\item Sticky skating - Regular, Figure 8, low
\item Weaves focusing on 180 degree turns 
\item Crossover-Carves
\item One foot lateral glides
\item 1 minute crossovers, 1 minute backwards crossovers, 1 minute anti-derby crossovers, 1 minute anti-derby backwards crossovers.  
\end{itemize}

Followed by any dynamic stretching and stopping. 

\subsection*{One Foot Blocking - 10 minutes}

Start with two foot blocking to find edges if needed.

Warming up edges with some contact. 
One skater driving, the other skater blocking on one foot. 

Extend with some more exotic drives for offence practice. 


\begin{itemize}
\item \hyperref[drill:drives/hip]{Hip drive} starting engaged
\item Hip drive from one to two feet 
\item \hyperref[drill:drives/shoulder]{Shoulder drive} 
\item \hyperref[drill:drives/low_shoulder]{Low shoulder drive}- with a focus on shoulder to butt cheek driving where possible. 
\item Shoulder drive into \hyperref[drill:drives/chest]{Chest drive}
\item \hyperref[drill:drives/shelf]{Shelf drive} starting engaged.
\end{itemize}



\subsection*{Supported Drive Drills - 10 minutes}

Driving using a two-wall or a three wall focus on getting a fast pocket or a hit-out.
Start with the jammer a foot or two behind the wall.
Combine with cutting your friend on the hit-out.


\subsection*{Joint Powerjam Drills}
This next section involves pairing strategies and powerjam scenarios.
Skaters should be divided between postive and negative teams, and be given one of the drills.

Positive:
\begin{itemize}
\item Offence including \hyperref[drill:drives/shelf]{chasing down}
\item Hold a goat and offence
\end{itemize}

Negative:
\begin{itemize}
\item \hyperref[drill:power_jam/running_away]{Running away}, or recycle at the front
\item I\hyperref[drill:power_jam/infinite_recycle]{Infinite recycle} at the back, pulling the jammer back to avoid offence.
\end{itemize}

\subsection*{Joint Power Jam Starts}

Having established what position we want to achieve in a powerjam, we now discuss how to start one. 

Positive power-jam team have the tempo, negative powerjam team need to both form and avoid being goated.

Suggestion for the negative team is to start centred but with a plan to take the jammer or pivot line.
If the pivot line, then the plan is to run, and recycle at the front, if the jammer line then the plan is for an immediate recycle and run back. 


\pagebreak


\section*{Drill Notes}

\input{../../drills/contact/solo_blocking/driving/hip}
\subsection*{Drill: Shoulder Drive}
\label{drill:drives/shoulder}

This drill requires pairs of skaters. 

Starting engaged, one skater is driving the other blocking, with both skaters starting in the middle of the track. 

The driving skater should attempt to use their shoulder to drive the other blocker.  
The skater being driven should attempt to hold their position (rather than simply skating away).


In this case, the shoulder can be used to partially lift the other blocker.
The idea is to set up in the armpit or ribs of the other blocker, and leverage from this position.
As friction requires that they're applying weight into the floor, this drive effectively negates the blockers ability to engage their edges with one leg.    


Practice in pairs, and rotate skaters where possible - driving is best practiced against a wide range of different opponents.
Also don't forget to practice driving to both the inside and outside.


% TODO: Figure

\subsection*{Drill: Low Shoulder Drive}
\label{drill:drives/shoulder_low}

This drill requires pairs of skaters. 

Starting engaged, one skater is driving the other blocking, with both skaters starting in the middle of the track. 

The driving skater should attempt to use their shoulder to drive the other blocker.  
The skater being driven should attempt to hold their position (rather than simply skating away).


In this case, the shoulder can be used to partially rotate the other blocker. 
The driving skater should set their shoulder up next to the butt cheek of the blocker. Using this as a fulcrum they can then apply pressure to rotate the blocker. 
Once the blocker is rotated, they've effectively had their feet turned, as in the hip drive, and can now be pushed more easily. 

For this secondary drive, try moving from the shoulder drive into a chest drive to corale the blocker off the track. 


Practice in pairs, and rotate skaters where possible - driving is best practiced against a wide range of different opponents.
Also don't forget to practice driving to both the inside and outside.



% TODO: Figure

\subsection*{Drill: Chest Drive}
\label{drill:drives/chest}

This drill requires pairs of skaters. 

Starting engaged, one skater is driving the other blocking, with both skaters starting in the middle of the track. 

The driving skater should attempt to use their chest to drive the other blocker.  
The skater being driven should attempt to hold their position (rather than simply skating away).

The idea of this drive is to push from a chest-on-chest position, using a combination of the direction of play, and the momentum from an initial, alternative drive to get the skater moving.     

This is best started from a shoulder drive, either low or regular, then ends in a chest drive. 
Keeping one leg down-track from the blocker prevents them from rotating in that direction, ideally getting them onto toe-stops further limits their ability to get out of the situation, which is helped by lifting them with the initial shoulder drive. 


Practice in pairs, and rotate skaters where possible - driving is best practiced against a wide range of different opponents.
Also don't forget to practice driving to both the inside and outside.


% TODO: Figure

\subsection*{Drill: Shelf Drive}
\label{drill:drives/shelf}

This drill requires pairs of skaters. 

Starting engaged, one skater is driving the other blocking, with both skaters starting in the middle of the track. 

The driving skater should attempt to shelf and drive the other blocker. 
The skater being driven should attempt to hold their position (rather than simply skating away).


This is a mix of the hip and shoulder drive.
Setting up with a leg placed laterally beneath the blocker, and a shoulder in the same position as the shoulder drive, the idea is to perform a lateral directly into the blocker. 
It is important to not engage the toe stop under the target blocker, as this will burn off the momentum of the drive.
Additionally it is important to start the engagement with either the hip, or the shoulder - starting with the thigh under the other blocker is a back block if any impact is achieved, and the purpose of the drive is to cause impact.


Using a `skateboard' lateral push, or a toe-stop push continue to drive the other blocker off the track.
The shoulder and thigh set up prevents the blocker from getting low - limiting the friction they can generate to resist the push.



Practice in pairs, and rotate skaters where possible - driving is best practiced against a wide range of different opponents.
Also don't forget to practice driving to both the inside and outside.



% TODO: Figure



\subsection*{Drill: Running Away (Negative)}
\label{drill:power_jam/running_away}

This drill requires a three wall, a jammer at at least one offence.
Optionally this may be extended to a four wall, and up to four offences.  
This drill may be paired with a Positive power-jam drill for the offensive skaters. 


Starting with the four-wall 8-10 feet down track from the offence with the jammer another 5 to 10 feet behind, the goal is to maintain pack while blocking the jammer, but keep the offence away from the wall.

The goal of the offence is to get their jammer out.
This may be through holding position to try to break pack, taking a goat and holding pack, or attempting a regular offence. 

The defence in this first drill should focus on running away while holding the inside line.
If the defence has the luxury of a fourth blocker, they may use one blocker as a bridge whose additional job is to communicate when the wall should move.
The engagement zone extends 20ft from the bridge, while the bridge may hold up to 10 feet ahead of the foremost offence.
For safer ranges, the bridge should hold at around 8 feet from the offence, while the wall should be 8-15 feet from the bridge.  


\begin{itemize}
\item If the offence runs up, the bridge should communicate and the wall move up to maintain distance from the jammer.
\item If the offence stops then bridge should hold position. The defensive wall should then seek to recycle the jammer back to the bridge and reset. 
\item Goats should use appropriate tactics (see a goating drill for that).
\end{itemize}




\subsection*{Drill: Infinite Recycle (Negative)}
\label{drill:power_jam/infinite_recycle}

This drill requires a three wall, a jammer at at least one offence.
Optionally this may be extended to a four wall, and up to four offences.  
This drill may be paired with a Positive power-jam drill for the offensive skaters. 


This drill starts with the jammer two feet behind the three wall, which is in turn 5-10 feet behind the other blockers.
A fourth blocker may position themselves freely, possibly running counter-offence, or may start as a goat.

The blocker's goal is to immediately recycle the jammer - taking themselves out if required so that another blocker may maintain position. 
If they do take themselves out then they are strongly encouraged to cut-one-friend on their way back in. 

The offence now serves to extend the engagement zone backwards, assisting in delaying the jammer.
The counter-offence skater acts to manage the pack, and communicates how far the run-back can go. 

As the jammer re-enters, the wall should recatch and pocket them - if offence comes back to assist then the jammer should be again hit out and recycled extending off the offence.

%TODO FIGURE

\subsection*{Drill: Goated at the Back (Negative)}
\label{sec:power_jam/goated}

The goat, the bunny, effectively this is a blocker that the other team has captured to maintain control of pack. 

Typically this will be at the back of pack, which is the case we will consider in this drill. 


Goat control depends on the pack definition and the engagement zone.



If the offence takes a goat (typically the bridge), the goat should attempt to escape. If the defensive wall is less than 10 feet away or recycling the jammer, they should aim to have themselves hit off track. If the defensive wall is more than 10 feet away, but still in the engagement zone, being hit off will trigger a no-pack and allow the jammer to escape. In this situation they should aim to push up. 


% TODO FIGURE

\subsection*{Drill: Chasing Down (Positive)}
\label{drill:power_jam/chasing_down}

This drill requires a three wall, a jammer at at least one offence.
Optionally this may be extended to a four wall, and up to four offences.  
This drill may be paired with a Negative power-jam drill for the defensive skaters. 


Starting with the four-wall 8-10 feet down track from the offence with the jammer another 5 to 10 feet behind, the goal is to maintain pack while blocking the jammer, but keep the offence away from the wall


The offensive skaters should start with a run-down and try to perform offence on the wall. 
It might be an idea to let the jammer hit first to fix the wall's position.


To achieve this run-down there are two main ideas:
\begin{itemize}
\item Taking an unguarded inside line - the path distance is shorter and so at the same speed the offensive skaters will be able to drive the opposing blocker to the slower outside of the track, allowing their jammer to pass. 
\item Out-pacing the wall on the outside - this requires that at least one blocker is faster than the wall they are chasing.
\end{itemize}

In both cases the secondary goal of the offence should be to take a goat without hitting them out, this will give them pack control.


%TODO FIGURE



\end{document}
